\documentclass[a4paper,12pt]{article}
\usepackage[T2A]{fontenc}
\usepackage[utf8]{inputenc}
\usepackage[russian]{babel}
\usepackage{amsmath,amssymb,mathtools}
\renewcommand{\familydefault}{\sfdefault}
\usepackage{icomma}
\usepackage{geometry}
\geometry{top=22mm,bottom=23mm,left=25mm,right=25mm}
\usepackage{microtype}
\usepackage{tikz}
\usetikzlibrary{arrows.meta,positioning,shapes.geometric,calc}
\usepackage{booktabs,tabularx}
\usepackage{enumitem}
\usepackage{setspace}
\usepackage{needspace}
\usepackage{fancyhdr}
\usepackage{xcolor}
\usepackage{hyperref}
\definecolor{accent}{HTML}{008080}
\definecolor{darkaccent}{HTML}{005C60}
\definecolor{textgray}{HTML}{30343A}
\definecolor{softgray}{HTML}{6E7478}
\hypersetup{colorlinks=true,linkcolor=darkaccent}
\onehalfspacing
\setlength{\parindent}{0pt}
\setlength{\parskip}{4pt}
\setlist[enumerate]{itemsep=7pt,topsep=5pt,leftmargin=6mm}
\setlist[itemize]{itemsep=2pt,topsep=3pt,leftmargin=6mm}
\pagestyle{fancy}
\fancyhf{}
\fancyhead[L]{\small\textcolor{softgray}{Степени, логарифмы и тригонометрия}}
\fancyhead[R]{\small\textcolor{softgray}{Ярослав Гаврилов}}
\fancyfoot[C]{\small\thepage}
\renewcommand{\headrulewidth}{0.35pt}
\renewcommand{\footrulewidth}{0pt}
\setlength{\headheight}{15pt}
\newcommand{\sect}[2]{%
  \Needspace{7\baselineskip}%
  \vspace{6pt}\noindent{\large\textcolor{darkaccent}{#1.\enspace #2}}\par\vspace{-1pt}%
  {\color{accent}\rule{\linewidth}{0.45pt}}\par\vspace{5pt}}
\newcommand{\idea}[1]{\textcolor{darkaccent}{\textbf{Ключевая идея.}} #1}
\newcommand{\warn}[1]{\textcolor{darkaccent}{\textbf{Проверка.}} #1}
\newcommand{\ex}[1]{\textcolor{darkaccent}{\textbf{Пример.}} #1}
\begin{document}

% Титульный лист
\thispagestyle{empty}
\begin{center}
\vspace*{25mm}
{\normalsize\textcolor{darkaccent}{МАТЕМАТИКА}}\par
\vspace{24mm}
{\LARGE Чек-лист}\par
\vspace{10mm}
{\Large\textcolor{accent}{Комплексное повторение:\\[3pt] степени, логарифмы, тригонометрия}}\par
\vspace{20mm}
{\normalsize Подготовка к ЕГЭ по профильной математике}\par
\vfill
{\large Ярослав Гаврилов}\par
\vspace{12mm}
{\normalsize 2 ноября 2026 г.}\par
\vspace{17mm}
{\small Лёвин Артём Александрович\\[3pt] эксклюзивно для Ярослава Гаврилова}\par
\vspace{18mm}
\begin{tikzpicture}[x=1cm,y=1cm]
\draw[accent,line width=0.75pt] (0,0) .. controls (1.2,0.55) and (2.6,-0.45) .. (4.0,0)
 .. controls (5.2,0.5) and (6.7,-0.5) .. (8.3,0);
\end{tikzpicture}
\vspace{7mm}
\end{center}
\clearpage

\sect{1}{Маршрут комплексного повторения}
Сегодняшние задачи требуют быстро распознавать тип выражения и выбирать одно полезное преобразование. Используйте краткий маршрут:
\begin{center}
\textcolor{darkaccent}{\textbf{основание степени}}
\quad$\longrightarrow$\quad
\textcolor{darkaccent}{\textbf{свойство логарифма}}
\quad$\longrightarrow$\quad
\textcolor{darkaccent}{\textbf{угол}}\par
\vspace{1pt}
\small Переключайтесь между приёмами по условию каждой задачи.
\end{center}
\begin{itemize}
\item \textbf{Степени:} привести обе части к одинаковому основанию.
\item \textbf{Логарифмы:} раскрыть произведение или степень; проверить ОДЗ.
\item \textbf{Тригонометрия:} свести углы к табличным и согласовать единицы.
\item \textbf{Прикладная задача:} решить неравенство, отобрать время на нужном отрезке, найти долю.
\end{itemize}

\sect{2}{Показательные уравнения}
\idea{При одинаковом положительном основании, отличном от единицы, можно приравнять показатели.}
\[
 a^{F(x)}=a^{G(x)},\quad a>0,\ a\ne1
 \quad\Longleftrightarrow\quad F(x)=G(x).
\]
\ex{Приведение к двойке:}
\[
 4^{3-x}=\sqrt{8}
 \quad\Longleftrightarrow\quad
 2^{6-2x}=2^{3/2}
 \quad\Longleftrightarrow\quad
 6-2x=\frac32.
\]
Отсюда $x=\dfrac94$. Помните: $(a^m)^n=a^{mn}$, $\sqrt a=a^{1/2}$ при $a\ge0$.

\sect{3}{Свойства логарифмов и перемена основания}
Для $a>0$, $a\ne1$, $b>0$, $c>0$:
\[
\begin{aligned}
\log_a(bc)&=\log_a b+\log_a c,\\
\log_a b^r&=r\log_a b\quad(b>0),\\
\log_a a^r&=r,\qquad \log_a1=0.
\end{aligned}
\]
\textbf{Взаимно обратные логарифмы:}
\[
 \log_a b=\frac{1}{\log_b a},\qquad a,b>0,\quad a\ne1,\quad b\ne1.
\]
\ex{Если $\log_b a=-2$, вычислите $\log_a(b^6\sqrt a)$:}
\[
\begin{aligned}
\log_a(b^6\sqrt a)&=6\log_a b+\tfrac12
 =\frac{6}{-2}+\frac12=\boxed{-\frac52}.
\end{aligned}
\]
\warn{Числа в аргументе логарифма обязаны быть положительными. Основание также положительно и не равно единице.}

\sect{4}{Логарифмические уравнения}
\idea{Сперва записываем ограничения на аргументы. Затем при одинаковом допустимом основании сравниваем аргументы и проверяем найденные корни.}
\ex{Решите уравнение:}
\[
 \log_3\left(\frac14-x\right)=\log_3(x^2-x).
\]
ОДЗ: $\dfrac14-x>0$ и $x^2-x>0$. Тогда
\[
 \frac14-x=x^2-x\quad\Longrightarrow\quad
 x^2=\frac14\quad\Longrightarrow\quad x=\pm\frac12.
\]
При $x=\tfrac12$ первое подлогарифмическое выражение отрицательно. При $x=-\tfrac12$ оба положительны. \textbf{Ответ:} $\boxed{-\tfrac12}$.

\textbf{Если переменная стоит в показателе под логарифмом,} полезно сразу вынести показатель:
\[
 \log_8\bigl(32^{5-x}\bigr)=4
 \ \Longrightarrow\ \frac53(5-x)=4
 \ \Longrightarrow\ \boxed{x=\frac{13}{5}}.
\]
Основание $8=2^3$, число $32=2^5$, поэтому $\log_8 32=\frac53$.

\clearpage
\sect{5}{Логарифмические графики: какие точки полезны}
Для функции
\[
 f(x)=k\log_a x,\qquad x>0,\quad a>0,\ a\ne1,
\]
точка $(1;0)$ принадлежит графику при любых допустимых $k$ и $a$, поскольку $\log_a1=0$. Поэтому она \textbf{не помогает определить параметры}.

\ex{Известно, что на графике лежат $(3;2)$ и $(9;4)$. Найдите $x$, если $f(x)=6$.}
\[
 f(3)=k\log_a3=2,
 \qquad f(9)=k\log_a(3^2)=2f(3)=4.
\]
Второе равенство повторяет информацию первого. Однако
\[
 f(27)=k\log_a(3^3)=3k\log_a3=3\cdot2=6.
\]
\textbf{Ответ:} $\boxed{x=27}$. При $k\ne0$ функция монотонна, поэтому нужное $x$ единственно.

\warn{Эти две точки не определяют $k$ и $a$ по отдельности. Для ответа достаточно свойства $\log_a(x^n)=n\log_a x$.}

\sect{6}{Логарифмы с разными основаниями}
Если в дроби встречаются логарифмы разных оснований, представьте основания и аргументы степенями:
\[
 \frac{\log_4 125}{\log_8 625}
 =\frac{3\log_{2^2}5}{4\log_{2^3}5}
 =\frac{(3/2)\log_2 5}{(4/3)\log_2 5}
 =\boxed{\frac98}.
\]
\idea{Степень аргумента выносится множителем, степень основания превращается в делитель: $\log_{a^m}b=\frac1m\log_a b$, если $m\ne0$ и основание допустимо.}

\sect{7}{Тригонометрия: единые углы}
\textbf{Двойной угол:}
\[
 \sin2\alpha=2\sin\alpha\cos\alpha.
\]
\textbf{Периодичность и чётность:}
\[
\begin{gathered}
\sin(\alpha+2\pi n)=\sin\alpha,\quad
\cos(\alpha+2\pi n)=\cos\alpha,\\
\tg(\alpha+\pi n)=\tg\alpha,\quad
\cos(-\alpha)=\cos\alpha,\\
\sin(-\alpha)=-\sin\alpha,\qquad
\tg(-\alpha)=-\tg\alpha.
\end{gathered}
\]
Здесь $n\in\mathbb Z$; у тангенса $\cos\alpha\ne0$.

\ex{Согласование углов в градусах:}
\[
 \frac{10\sin85^\circ\cos85^\circ}{\sin190^\circ}
 =\frac{5\sin170^\circ}{\sin190^\circ}
 =\frac{-5\sin190^\circ}{\sin190^\circ}
 =\boxed{-5}.
\]
Использованы $170^\circ=360^\circ-190^\circ$ и
$\sin(360^\circ-\alpha)=-\sin\alpha$.

\warn{Не смешивайте градусы и радианы. Пишите $^\circ$ в каждом выражении, где угол задан в градусах.}

\sect{8}{Приведение больших углов к табличным}
Если угол выражен в радианах, сначала вычтите целое число периодов, затем воспользуйтесь табличным углом.
\[
 \sin\frac{23\pi}{6}=\sin\left(4\pi-\frac{\pi}{6}\right)=-\frac12,
\]
\[
 \cos\left(-\frac{11\pi}{4}\right)
 =\cos\frac{11\pi}{4}=\cos\frac{3\pi}{4}=-\frac{\sqrt2}{2},
\]
\[
 \tg\frac{8\pi}{3}=\tg\frac{2\pi}{3}=-\sqrt3.
\]
Поэтому
\[
 \frac{6\sin(23\pi/6)\cos(-11\pi/4)}{\tg(8\pi/3)}
 =\boxed{-\frac{\sqrt6}{2}}.
\]
\idea{Для синуса и косинуса период $2\pi$, для тангенса --- $\pi$. Проверяйте знак по четверти единичной окружности.}

\sect{9}{Прикладная задача: доля времени}
\textbf{Формулировка.} Скорость (в см/с) на отрезке наблюдения $0\le t\le3{,}5$ задана формулой
\[
 v(t)=2\sqrt2\sin\left(\frac{\pi t}{6}\right).
\]
Какую долю первых $3{,}5$ с скорость \textbf{не превышала} $2$ см/с?

\textbf{Шаг 1.} Переводим слова в неравенство:
\[
 2\sqrt2\sin\frac{\pi t}{6}\le2
 \quad\Longleftrightarrow\quad
 \sin\frac{\pi t}{6}\le\frac{\sqrt2}{2}.
\]
\Needspace{8\baselineskip}
\textbf{Шаг 2.} На окружности подходят углы, для которых синус не больше $\sqrt2/2$. Один из периодических отрезков:
\[
 -\frac{5\pi}{4}+2\pi n\ \le\ \frac{\pi t}{6}\ \le\ \frac\pi4+2\pi n,
 \qquad n\in\mathbb Z.
\]
На промежутке $0\le t\le3{,}5$ достаточно рассмотреть $n=0$: $-7{,}5\le t\le1{,}5$. Его пересечение с отрезком наблюдения --- $[0;1{,}5]$.

\textbf{Шаг 3.} Продолжительность подходящего времени равна $1{,}5$ с, а полный интервал наблюдения равен $3{,}5$ с. Доля времени:
\[
 \frac{1{,}5}{3{,}5}=\frac37=0{,}428571\ldots\ \approx\ \boxed{0{,}43}.
\]
\warn{При другом интервале наблюдения может потребоваться несколько периодических отрезков. Берите их пересечения с $[0;T]$ и складывайте длины.}

\sect{10}{Памятка по типичным ошибкам}
\begin{itemize}
\item \textbf{Степени:} при $(a^m)^n$ показатели перемножаются, а при $a^m a^n$ складываются.
\item \textbf{Логарифмы:} запись $\log_a(u+v)$ нельзя разделять на сумму логарифмов.
\item \textbf{ОДЗ:} после решения логарифмического уравнения всегда проверяйте найденные корни.
\item \textbf{График:} точка $(1;0)$ для $k\log_a x$ не определяет коэффициенты.
\item \textbf{Тригонометрия:} период, четверть и знак обязательны; для тангенса учитывайте запрет деления на нуль.
\item \textbf{Прикладная задача:} доля $=$ подходящее время $/$ всё время наблюдения; единицы измерения должны согласовываться.
\end{itemize}

\clearpage
\thispagestyle{fancy}
\begin{center}
{\Large\textcolor{darkaccent}{Алгоритм: тригонометрическое условие на время}}\par
\vspace{11mm}
\end{center}
\tikzset{
 myblock/.style={rectangle,rounded corners=2pt,draw=accent,line width=0.65pt,align=center,text width=100mm,minimum height=11mm,inner sep=4pt,font=\small},
 mydecision/.style={diamond,aspect=3.1,draw=accent,line width=0.65pt,align=center,text width=64mm,inner sep=2pt,font=\small},
 myoval/.style={ellipse,draw=darkaccent,line width=0.75pt,align=center,minimum width=55mm,minimum height=10mm,font=\small},
 myarrow/.style={-{Latex[length=2.6mm]},draw=darkaccent,line width=0.75pt}
}
\begin{center}
\begin{tikzpicture}[node distance=9mm]
\node[myoval] (start) {Дано $v(t)$, наблюдаем $0\le t\le T$};
\node[myblock,below=of start] (ineq) {Перевести условие «не превышает» в $v(t)\le v_0$};
\node[myblock,below=of ineq] (unit) {Решить тригонометрическое неравенство для угла $u$};
\node[myblock,below=of unit] (periods) {Учесть периодические сдвиги $2\pi n$ или $\pi n$};
\node[myblock,below=of periods] (time) {Перейти от углов $u$ к промежуткам времени $t$};
\node[myblock,below=of time] (intersection) {Пересечь все подходящие промежутки с $[0;T]$};
\node[mydecision,below=13mm of intersection] (has) {Есть пересечения?};
\node[myblock,below=14mm of has] (sum) {Да: сложить их длины, получить время $L$};
\node[myblock,right=14mm of has,text width=25mm,minimum height=13mm] (zero) {Нет: $L=0$};
\node[myoval,below=10mm of sum] (finish) {Доля $=L/T$};
\draw[myarrow] (start)--(ineq);
\draw[myarrow] (ineq)--(unit);
\draw[myarrow] (unit)--(periods);
\draw[myarrow] (periods)--(time);
\draw[myarrow] (time)--(intersection);
\draw[myarrow] (intersection)--(has);
\draw[myarrow] (has)--node[right,font=\scriptsize,pos=.38]{да}(sum);
\draw[myarrow] (has.east)--node[above,font=\scriptsize]{нет}(zero.west);
\draw[myarrow] (sum)--(finish);
\draw[myarrow] (zero.south) |- (finish.east);
\end{tikzpicture}
\end{center}
\clearpage

\begin{center}
{\Large\textcolor{darkaccent}{Тренировочный блок}}\par
{\small Путь А · 10 заданий · от простого к сложному}
\end{center}
\vspace{3mm}
Решайте последовательно; в логарифмических уравнениях фиксируйте ОДЗ, а в тригонометрии следите за знаками и единицами угла.
\begin{enumerate}[label=\textcolor{darkaccent}{\textbf{\arabic*.}},itemsep=8pt]
\item Решите уравнение: $4^{2-x}=\sqrt8$.
\item Решите уравнение: $\log_{27}\left(9^{x+1}\right)=2$.
\item Известно, что $\log_b a=-3$. Вычислите $\log_a(b^4\sqrt a)$.
\item Решите уравнение:
$\displaystyle\log_5\left(\frac19-x\right)=\log_5(x^2-x)$.
\item Вычислите:
$\displaystyle\frac{\log_9 125}{\log_{27}625}$.
\item Вычислите:
$\displaystyle\frac{6\sin(19\pi/6)\cos(-7\pi/4)}{\tg(5\pi/4)}$.
\item Вычислите:
$\displaystyle\frac{12\sin65^\circ\cos65^\circ}{\sin230^\circ}$.
\item График $f(x)=k\log_a x$ проходит через точки $(1;0)$, $(2;3)$ и $(4;6)$. Найдите $x$, при котором $f(x)=9$.
\item Скорость задана формулой $\displaystyle v(t)=2\sqrt2\sin\frac{\pi t}{8}$ (см/с). Какую долю первых $4$ секунд скорость не превышает $2$ см/с? Ответ выразите десятичной дробью.
\item Скорость задана формулой $\displaystyle v(t)=2\sqrt2\sin\frac{\pi t}{6}$ (см/с). Какую долю первых $5$ секунд скорость не превышает $2$ см/с? Ответ выразите десятичной дробью.
\end{enumerate}
\vspace{4mm}
\textcolor{darkaccent}{\textbf{Самопроверка перед сдачей:}}\enspace проверены ОДЗ, знаки тригонометрических значений, интервалы времени и округление.
\clearpage

\begin{center}
{\Large\textcolor{darkaccent}{Ответы}}\par
\vspace{5mm}
\end{center}
\renewcommand{\arraystretch}{1.45}
\begin{center}
\begin{tabularx}{.91\linewidth}{@{}>{\bfseries}p{13mm}X@{}}
\toprule
№ & Ответ \\
\midrule
1 & $x=\dfrac54$ \\
2 & $x=2$ \\
3 & $-\dfrac56$ \\
4 & $x=-\dfrac13$ \\
5 & $\dfrac98$ \\
6 & $-\dfrac{3\sqrt2}{2}$ \\
7 & $-6$ \\
8 & $x=8$ \\
9 & $0{,}5$ \\
10 & $0{,}4$ \\
\bottomrule
\end{tabularx}
\end{center}
\vspace{12mm}
\sect{11}{Итоговая памятка}
\begin{tabularx}{\linewidth}{@{}p{40mm}X@{}}
\toprule
\textbf{Сигнал в задаче} & \textbf{Первое полезное действие}\\
\midrule
Составные основания & Представьте их степенями одного числа.\\
Сумма логарифмов & Используйте свойство логарифма произведения.\\
Переставлены основание и аргумент & Примените формулу взаимно обратных логарифмов.\\
Одинаковые логарифмы & Запишите ОДЗ и приравняйте аргументы.\\
График $k\log_a x$ & Ищите связь между $x$, $x^2$, $x^3$.\\
Разные углы & Используйте двойной угол, периодичность и формулы приведения.\\
Доля времени & Пересеките найденные интервалы с отрезком наблюдения.\\
\bottomrule
\end{tabularx}
\end{document}
